\section*{Глава 2. Проектирование веб-приложения}
\addcontentsline{toc}{section}{Глава 2. Проектирование веб-приложения}
\stepcounter{section}


\subsection{Анализ целевой аудитории}

Для успешной разработки веб-приложения по ведению семейного бюджета необходимо чётко понимать, для кого оно создаётся, какие задачи будут решать пользователи и с какими ограничениями (техническими, финансовыми, временными) они сталкиваются в повседневной жизни.

На основе потребностей, технической подготовки и мотивации к ведению учёта была выделена одна, самая многочисленная и проблемная категории пользователей: семьи с совместным (полностью или частично) бюджетом.

\subsubsection{Проблематика}

Семьи с совместным бюджетом сталкиваются с рядом типичных проблем:
\begin{itemize}
    \item Отсутствие финансовой прозрачности. Часто только один из партнёров имеет полное представление о доходах и расходах, что порождает недоверие.
    \item Споры о тратах и невозможность объективной оценки. При отсутствии единой записной книжки невозможно объективно сравнить расходы разных членов семьи.
    \item Сложность планирования крупных покупок. При обдумывании больших трат необходимо иметь актуальную картину по доходам и обязательным расходам.
    \item Необходимость контролировать разные типы расходов. Семейный бюджет неоднороден: есть общие траты и личные траты. Без разграничения можно незаметить лишние расходы.
    \item Разная мотивация к ведению учёта. Один партнёр может хотеть детальной аналитики, а второй лишь изредка вносить расходы. Инструмент должен учитывать оба типа поведения.
\end{itemize}

\subsubsection{Потребности ЦА}

На основе выявленных проблем были сформулированы следующие потребности, которые должно закрыть разрабатываемое приложение:

\begin{table}[h]
    \centering
    \caption{Потребности целевой аудитории}
    \label{tab:target-audience-needs}

    {\fontsize{12}{14}\selectfont
	\setlength{\tabcolsep}{6pt}

    \begin{tabular}{|p{4cm}|p{8cm}|p{3cm}|}
        \hline
        \textbf{Потребность} & \textbf{Описание} & \textbf{Задача из ТЗ} \\
        \hline
        Единое хранилище & Все транзакции семьи в одном месте, доступ с любого устройства & 3.3 (CRUD) \\
        \hline
        Разделение по ролям & Администратор, родитель, ребенок & 3.2 (авторизация с ролями) \\
        \hline
        Категоризация трат & Общие и личные категории, возможность создания/редактирования & 3.3 (категории) \\
        \hline
        Агрегация по участникам & Автоматический подсчёт расходов каждого члена семьи за период & 3.5 (аналитика) \\
        \hline
        Базовые отчёты & По периодам, категориям, участникам & 3.5 (статистика) \\
        \hline
        Привязка чеков & Загрузка фото/скана чека к транзакции & 3.4 (файлы) \\
        \hline
        Наглядная визуализация & Диаграммы и графики для быстрого понимания структуры расходов & 3.5 (Chart.js) \\
        \hline

    \end{tabular}
    }
\end{table}

\subsubsection{Портрет типичного представителя категории}

\begin{itemize}
    \item Семейное положение: состоит в браке или зарегистрированном партнёрстве.
    \item Состав семьи: 2–4 человека (родители + 1–2 детей).
    \item Используемые устройства: смартфон, ноутбук или ПК.
\end{itemize}

\subsubsection{Ключевые функции приложения}

На основе анализа проблем и потребностей выделен основной набор функций:

\begin{table}[h]
    \centering
    \caption{Основные функции}
    \label{tab:main-functions}

    {\fontsize{12}{14}\selectfont
	\setlength{\tabcolsep}{6pt}

    \begin{tabular}{|p{8cm}|p{4cm}|p{3cm}|}
        \hline
        \textbf{Функция} & \textbf{Важность} & \textbf{Реализация} \\
        \hline
        Аутентификация и разделение ролей & Критическая & Задача 3.2 \\
        \hline
        CRUD для транзакций & Критическая & Задача 3.3 \\
        \hline
        Фильтрация по дате, категории, участнику & Высокая & Задача 3.3 \\
        \hline
        Управление категориями (CRUD) & Высокая & Задача 3.3 \\
        \hline
        Управление бюджетами/лимитами & Высокая & Задача 3.3 \\
        \hline
        Визуализация (диаграммы) & Высокая & Задача 3.5 \\
        \hline
        Загрузка чеков & Средняя & Задача 3.4 \\
        \hline
        Экспорт отчёта & Средняя & Задача 3.4 \\
        \hline

    \end{tabular}
    }
\end{table}

\subsection{Описание функциональности приложения}


\subsection{Проектирование модели данных}


\subsection{Разработка макетов страниц}

